\section*{अध्याय \ref{ch:b1:topology} --- वास्तविक रेखा की सांस्थिति}

\begin{solution}{exo:b1:topology:1}
$\intoo{0}{1} \cup \intoo{2}{3}$: \omterm{def:b1:topology:open}{विवृत} (\omterm{def:b1:topology:open}{विवृत} समुच्चयों का सम्मिलन), \omterm{def:b1:topology:closed}{संवृत}
नहीं ($\frac 1n \to 0$ बाहर)।

$\intco{0}{1}$: कुछ भी नहीं। \omterm{def:b1:topology:open}{विवृत} नहीं ($0$ के चारों ओर कोई \omterm{prop:b1:reals:intervals}{अंतराल} भीतर
नहीं); \omterm{def:b1:topology:closed}{संवृत} नहीं ($1 - \frac1n \to 1 \notin$ \omterm{def:b1:logic:sets}{समुच्चय})।

$\{0\} \cup \intcc{1}{2}$: \omterm{def:b1:topology:closed}{संवृत} (\omterm{def:b1:topology:closed}{संवृत} समुच्चयों का परिमित सम्मिलन), \omterm{def:b1:topology:open}{विवृत}
नहीं ($0$ पर विफल)।

$\R \setminus \Z$: \omterm{def:b1:topology:open}{विवृत} ($\Z$ \omterm{def:b1:topology:closed}{संवृत} है), \omterm{def:b1:topology:closed}{संवृत} नहीं: अनुक्रम $\bigl(\frac
1n\bigr)$ उसमें है, पर उसकी सीमा $0$ $\Z$ की सदस्य है, अर्थात् \omterm{def:b1:logic:sets}{समुच्चय} से
भाग जाती है।

$\Q \cap \intoo{0}{1}$: कुछ भी नहीं। \omterm{def:b1:topology:open}{विवृत} नहीं: किसी परिमेय संख्या के चारों
ओर हर \omterm{prop:b1:reals:intervals}{अंतराल} में अपरिमेय संख्याएँ हैं। \omterm{def:b1:topology:closed}{संवृत} नहीं: उसमें ऐसे अनुक्रम हैं जो
अपरिमेय $\frac{\sqrt 2}{2}$ की ओर जाते हैं (\omterm{def:b1:topology:dense}{सघनता})।
\end{solution}

\begin{solution}{exo:b1:topology:2}
$A = \intoc{0}{1} \cup \{2\}$: $\mathring A = \intoo{0}{1}$,
$\overline A = \intcc{0}{1} \cup \{2\}$, $\partial A = \{0, 1, 2\}$.

$A = \R \setminus \Q$: $\mathring A = \emptyset$ (हर \omterm{prop:b1:reals:intervals}{अंतराल} में परिमेय
संख्याएँ हैं), $\overline A = \R$ (अपरिमेय संख्याओं की \omterm{def:b1:topology:dense}{सघनता}), $\partial A =
\R$।

$A = \bigl\{\frac{(-1)^n n}{n+1}\bigr\}$: सम पद $1$ की ओर जाते हैं और विषम
पद $-1$ की ओर, और इनमें से कोई $A$ का सदस्य नहीं है। $\mathring A =
\emptyset$ (वियुक्त बिंदु), $\overline A = A \cup \{-1, 1\}$, $\partial A =
\overline A$।
\end{solution}

\begin{solution}{exo:b1:topology:3}
\emph{पूरक।} $F = \{a_1 < a_2 < \dots < a_k\}$: पूरक \omterm{def:b1:topology:open}{विवृत} अंतरालों
$\intoo{-\infty}{a_1}$, $\intoo{a_i}{a_{i+1}}$, $\intoo{a_k}{+\infty}$ का
सम्मिलन है --- जो \cref{prop:b1:topology:openstable} से \omterm{def:b1:topology:open}{विवृत} है।

\emph{अनुक्रम।} मान लीजिए $u_n \in F$, $u_n \to \ell$। $\varepsilon =
\min\{\abs{a_i - a_j} : i \neq j\}/2 > 0$ के साथ (अथवा $F$ के एकल \omterm{def:b1:logic:sets}{समुच्चय}
होने पर कोई भी $\varepsilon$): किसी कोटि के आगे सभी पद $\ell$ के
$\varepsilon$ के भीतर हैं, अतः एक-दूसरे के $2\varepsilon$ के भीतर, जो उन्हें
उस कोटि से आगे एक ही $a_i$ होने पर बाध्य कर देता है; और तब $\ell = a_i \in
F$।
\end{solution}

\begin{solution}{exo:b1:topology:4}
$\subseteq$: $\overline A \cup \overline B$ \omterm{def:b1:topology:closed}{संवृत} है (परिमित सम्मिलन) और
उसमें $A \cup B$ समाया है, अतः उसमें \emph{सबसे छोटा} \omterm{def:b1:topology:closed}{संवृत} अधिसमुच्चय
$\overline{A \cup B}$ भी समाया है। $\supseteq$: $A \subseteq A \cup B$
$\overline A \subseteq \overline{A \cup B}$ देता है (\omterm{def:b1:topology:closure}{संवरक} एकदिष्ट है: $A$
के अभिलग्न बिंदु बड़े \omterm{def:b1:logic:sets}{समुच्चय} से भी अभिलग्न होते हैं), और $B$ के लिए भी वैसे
ही।

प्रतिच्छेदन के लिए प्रति-उदाहरण: $A = \intoo{0}{1}$, $B = \intoo{1}{2}$:
$\overline{A \cap B} = \overline\emptyset = \emptyset$, पर $\overline A \cap
\overline B = \{1\}$।
\end{solution}

\begin{solution}{exo:b1:topology:5}
अनुक्रमीय अभिलक्षण (\cref{thm:b1:topology:seqclosed}) का प्रयोग कीजिए। मान
लीजिए $S = \{u_n\} \cup \{\ell\}$ और $S$ में $v_k \to m$ वाला कोई अनुक्रम
$(v_k)$ है; दिखाइए $m \in S$। दो स्थितियाँ। यदि $(v_k)$ कोई मान $v \in S$
अनंत बार लेता है, तो अचर \omterm{def:b1:seq:subsequence}{उपानुक्रम} $m = v \in S$ दे देता है। अन्यथा हर मान
परिमित बार लिया जाता है; विशेष रूप से प्रत्येक $n$ के लिए पद $u_n$ परिमित
बार आता है, और $\ell$ भी। तब प्रत्येक $N$ के लिए $v_k \in \{u_0, \dots, u_N,
\ell\}$ वाले सूचकांक $k$ परिमित कितने हैं: शेष $v_k$ ऐसे पद $u_n$ हैं जिनके
लिए $n > N$। $\varepsilon > 0$ दिया हो, तो ऐसा $N$ चुनिए कि $n > N$ के लिए
$\abs{u_n - \ell} \leq \varepsilon$: तब परिमित कितने को छोड़कर सभी $v_k$
$\abs{v_k - \ell} \leq \varepsilon$ संतुष्ट करते हैं। अतः $v_k \to \ell$,
इसलिए $m = \ell \in S$।
\end{solution}

\begin{solution}{exo:b1:topology:6}
$U + A = \bigcup_{a \in A} (U + a)$, और प्रत्येक स्थानांतरण $U + a$ \omterm{def:b1:topology:open}{विवृत} है
(अंतराल-प्रमाणपत्रों का स्थानांतरण)। \omterm{def:b1:topology:open}{विवृत} समुच्चयों का सम्मिलन \omterm{def:b1:topology:open}{विवृत} होता
है (\cref{prop:b1:topology:openstable})।

\omterm{def:b1:topology:closed}{संवृत समुच्चय}: $\Z$ और $\sqrt 2\,\Z$ \omterm{def:b1:topology:closed}{संवृत} हैं (जैसे $\alpha\Z$: अभिसारी
अनुक्रम अंततः अचर होते हैं, \cref{ex:b1:topology:closed} देखिए)। उनका योग
$\Z + \sqrt 2\,\Z$ $\R$ में \omterm{def:b1:topology:dense}{सघन} है (\cref{exo:b1:reals:9}) पर वह $\R$ नहीं
है (वह गणनीय है; अथवा सरलता से: $\frac{\sqrt 2}{2} \notin \Z + \sqrt2\Z$,
अन्यथा $\sqrt 2$ परिमेय होता --- $\frac{\sqrt2}{2} = m + n\sqrt 2$ लिखने पर
$(2n - 1)\sqrt 2 = -2m$ अनिवार्य हो जाता, अतः $n = \frac12$ के बिना $\sqrt 2
\in \Q$, जो असंभव है)। कोई \omterm{def:b1:topology:dense}{सघन} उचित उपसमुच्चय \omterm{def:b1:topology:closed}{संवृत} नहीं होता: उसका \omterm{def:b1:topology:closure}{संवरक}
स्वयं $\R \neq$ है।
\end{solution}

\begin{solution}{exo:b1:topology:7}
$\Z$: $n$ का \omterm{def:b1:topology:open}{प्रतिवेश} $\intoo{n - \frac12}{n + \frac12}$ $\Z$ से केवल $n$ पर
मिलता है।

$A = \{\frac1n : n \in \N^*\}$: $\frac 1n$ के चारों ओर त्रिज्या $\frac{1}{n}
- \frac{1}{n+1} = \frac{1}{n(n+1)}$ (मान लीजिए आधी) का \omterm{prop:b1:reals:intervals}{अंतराल} उसे उसके
पड़ोसियों से अलग कर देता है --- अतः सभी बिंदु वियुक्त। फिर भी $0 \in
\overline A \setminus A$: वियुक्त बिंदु बाहरी अभिलग्न बिंदुओं को नहीं रोकते।

यदि $A$ के सभी बिंदु वियुक्त हैं: तो $A$ का कोई बिंदु आंतरिक नहीं है,
क्योंकि आंतरिक बिंदु के चारों ओर $A$-पड़ोसियों का पूरा \omterm{prop:b1:reals:intervals}{अंतराल} होता है (और
\omterm{prop:b1:reals:intervals}{अंतराल} अनंत होता है), जो वियुक्तता का विरोध करता है। अतः $\mathring A =
\emptyset$।
\end{solution}

\begin{solution}{exo:b1:topology:8}
यदि $A \subseteq \intcc{-M}{M}$, तो \omterm{def:b1:topology:closed}{संवृत समुच्चय} $\intcc{-M}{M}$ में $A$
समाया है, अतः उसमें $\overline A$ भी समाया है (सबसे छोटा \omterm{def:b1:topology:closed}{संवृत} अधिसमुच्चय):
अर्थात् $\overline A$ परिबद्ध है।

मान लीजिए $s = \sup A$ (परिमित)। चूँकि $A \subseteq \overline A$, $\sup
\overline A \geq s$। विलोमतः $\overline A \subseteq \intoc{-\infty}{s}$: वह
अर्धरेखा \omterm{def:b1:topology:closed}{संवृत} है और उसमें $A$ समाया है; अतः $\overline A$ का प्रत्येक अवयव
$\leq s$ है, जिससे $\sup\overline A \leq s$ मिलता है। समता।

$s \in \overline A$: $\varepsilon$-अभिलक्षण (\cref{prop:b1:reals:epsilon})
से प्रत्येक \omterm{prop:b1:reals:intervals}{अंतराल} $\intoo{s - \varepsilon}{s + \varepsilon}$ में $A$ का कोई
अवयव है: अतः $s$ अभिलग्न है।
\end{solution}

\begin{solution}{exo:b1:topology:9}
मान लीजिए $A$ \omterm{def:b1:topology:open}{विवृत} और \omterm{def:b1:topology:closed}{संवृत} है, जहाँ $a \in A$ और $b \in \R \setminus A$;
व्यापकता खोए बिना $a < b$। \omterm{def:b1:logic:sets}{समुच्चय} $B = A \cap \intcc{a}{b}$ अरिक्त ($a$) और
परिबद्ध है: मान लीजिए $s = \sup B$। \cref{exo:b1:topology:8} से $s \in
\overline B \subseteq \overline A = A$ ($A$ \omterm{def:b1:topology:closed}{संवृत})। ध्यान दीजिए $s \leq b$,
और चूँकि $b \notin A$: $s < b$। अब $A$ \omterm{def:b1:topology:open}{विवृत} है: कोई \omterm{prop:b1:reals:intervals}{अंतराल} $\intoo{s - r}{s
+ r}$ $A$ में है, और हम $r < b - s$ ले सकते हैं। तब $s + \frac{r}{2}$ $A
\cap \intcc{a}{b} = B$ का सदस्य है और $s$ से बड़ा है --- जो $s = \sup B$ का
विरोध करता है। अतः ऐसा कोई युग्म $(a, b)$ है ही नहीं: $A$, $\R \setminus A$
में से एक रिक्त है।
\end{solution}

\begin{solution}{exo:b1:topology:10}
$I_x$ ऐसे \omterm{def:b1:topology:open}{विवृत} अंतरालों का सम्मिलन है जिन सबमें $x$ है: अतः वह \omterm{def:b1:topology:open}{विवृत} है, और
उत्तल होने से \omterm{prop:b1:reals:intervals}{अंतराल} भी --- यदि $u, v \in I_x$ के साथ $u < z < v$, तो $u$ और
$v$ $U$ के \omterm{def:b1:topology:open}{विवृत} उपअंतरालों $J_u \ni x$, $J_v \ni x$ में हैं, और $J_u \cup
J_v$ ऐसा \omterm{prop:b1:reals:intervals}{अंतराल} है (दोनों में $x$ है) जो $U$ के भीतर है और $z$ को समेटता है;
अतः $z \in I_x$ (\cref{prop:b1:reals:intervals})।

यदि $I_x \cap I_y \neq \emptyset$: तो $I_x \cup I_y$ $U$ में समाया ऐसा \omterm{def:b1:topology:open}{विवृत}
\omterm{prop:b1:reals:intervals}{अंतराल} है (उत्तल: दो अतिव्यापी \omterm{prop:b1:reals:intervals}{अंतराल}) जो $x$ और $y$ दोनों को समेटता है, अतः
दोनों की अधिकतमता से $I_x \cup I_y \subseteq I_x$ और $\subseteq I_y$: $I_x =
I_y$।

अतः $U$ भिन्न समुच्चयों $I_x$ का असंयुक्त सम्मिलन है (प्रत्येक $x \in U$
अपने ही $I_x$ में है)। गणनीयता: इस परिवार के प्रत्येक अरिक्त \omterm{def:b1:topology:open}{विवृत} \omterm{prop:b1:reals:intervals}{अंतराल}
$I$ में कोई परिमेय संख्या $q_I$ है (\cref{thm:b1:reals:density}), और भिन्न
असंयुक्त अंतरालों को भिन्न परिमेय संख्याएँ मिलती हैं: अतः यह परिवार $\Q$ में
\omterm{def:b1:logic:inj}{एकैकी} रूप से समा जाता है, जो गणनीय है (वह पूर्णांक युग्मों से सूचीबद्ध है)।
अतः अधिक से अधिक गणनीय कितने \omterm{prop:b1:reals:intervals}{अंतराल}।
\end{solution}

\begin{solution}{exo:b1:topology:11}
$\overline A = A \cup A'$। ($\supseteq$) सदा $A \subseteq \overline A$; और
यदि $x \in A'$, तो $x$ का प्रत्येक \omterm{def:b1:topology:open}{प्रतिवेश} $A \setminus \{x\} \subseteq A$
से मिलता है, अतः $x$ अभिलग्न है। ($\subseteq$) मान लीजिए $x \in \overline
A$। यदि $x \in A$, तो काम पूरा। यदि $x \notin A$, तो $x$ का प्रत्येक
\omterm{def:b1:topology:open}{प्रतिवेश} $A = A \setminus \{x\}$ से मिलता है: $x \in A'$।

फलस्वरूप $A$ \omterm{def:b1:topology:closed}{संवृत} $\iff$ $A = \overline A = A \cup A'$ $\iff$ $A' \subseteq
A$।

$A = \{\frac 1n\}$: $0$ एक संचय बिंदु है ($\frac 1n \to 0$, पद $\neq 0$);
प्रत्येक $\frac 1n$ वियुक्त है (\cref{exo:b1:topology:7}), अतः $A'$ में
नहीं; और $x \notin A \cup \{0\}$ वाले बिंदु के चारों ओर पूरा \omterm{prop:b1:reals:intervals}{अंतराल} $A$ से
बचता है ($A \cup \{0\}$ में $x$ के दोनों पड़ोसियों के बीच, या $1$ से आगे)।
अतः $A' = \{0\}$।

$\Z' = \emptyset$: प्रत्येक पूर्णांक वियुक्त है, और प्रत्येक अपूर्णांक का
कोई \omterm{def:b1:topology:open}{प्रतिवेश} $\R \setminus \Z$ के भीतर है।

$\Q' = \R$: किसी भी वास्तविक संख्या के चारों ओर हर \omterm{prop:b1:reals:intervals}{अंतराल} में अनंत कितनी
परिमेय संख्याएँ हैं (\cref{thm:b1:reals:density}), और विशेष रूप से केंद्र से
भिन्न कोई एक।
\end{solution}

\begin{solution}{exo:b1:topology:12}
\begin{enumerate}
  \item प्रत्येक $a \in A$ के लिए: $\abs{x - a} \leq \abs{x - y} + \abs{y -
  a}$, अतः $d_A(x) \leq \abs{x - y} + \abs{y - a}$; $a$ पर \omterm{def:b1:reals:bounds}{निम्नतम} लेने पर:
  $d_A(x) \leq \abs{x - y} + d_A(y)$। $x$ और $y$ की अदला-बदली दूसरी असमिका
  देती है: $\abs{d_A(x) - d_A(y)} \leq \abs{x - y}$।
  \item $d_A(x) = 0$ $\iff$ प्रत्येक $\varepsilon > 0$ के लिए ऐसा $a \in A$
  है कि $\abs{x - a} < \varepsilon$ $\iff$ $x$ के चारों ओर प्रत्येक \omterm{prop:b1:reals:intervals}{अंतराल}
  $A$ से मिलता है $\iff$ $x \in \overline A$। यदि $F$ \omterm{def:b1:topology:closed}{संवृत} है और $x \notin
  F = \overline F$, तो $d_F(x) \neq 0$, अर्थात् $d_F(x) > 0$।
  \item यदि $d_F(x) < \varepsilon$, तो $r = \varepsilon - d_F(x) > 0$ रखिए:
  $\abs{y - x} < r$ के लिए भाग (1) $d_F(y) \leq d_F(x) + \abs{x - y} <
  \varepsilon$ देता है: अर्थात् \omterm{prop:b1:reals:intervals}{अंतराल} $\intoo{x - r}{x + r}$
  $V_\varepsilon$ में है, जो इसलिए \omterm{def:b1:topology:open}{विवृत} है; और उसमें $F$ है, क्योंकि वहाँ
  $d_F = 0$। अंत में $x \in \bigcap_{\varepsilon>0} V_\varepsilon$ $\iff$
  सभी $\varepsilon$ के लिए $d_F(x) < \varepsilon$ $\iff$ $d_F(x) = 0$ $\iff$
  $x \in \overline F = F$। चूँकि $\bigcap_{\varepsilon > 0} V_\varepsilon =
  \bigcap_{n \geq 1} V_{1/n}$, अतः प्रत्येक \omterm{def:b1:topology:closed}{संवृत समुच्चय} \omterm{def:b1:topology:open}{विवृत} समुच्चयों का
  गणनीय प्रतिच्छेदन है।
\end{enumerate}
\end{solution}

\begin{solution}{pb:b1:topology:1}
\textbf{1.} $C_1 = \intcc{0}{\frac13} \cup \intcc{\frac23}{1}$ और
\[
C_2 = \intcc{0}{\tfrac19} \cup \intcc{\tfrac29}{\tfrac13}
\cup \intcc{\tfrac23}{\tfrac79} \cup \intcc{\tfrac89}{1} .
\]
आगमन: यदि $C_n$ लंबाई $3^{-n}$ के $2^n$ \omterm{def:b1:topology:closed}{संवृत} खंडों का असंयुक्त सम्मिलन है,
तो हर एक का \omterm{def:b1:topology:open}{विवृत} मध्य-तिहाई हटाने से प्रति जनक लंबाई $3^{-n-1}$ के दो \omterm{def:b1:topology:closed}{संवृत}
खंड बच जाते हैं: $2^{n+1}$ खंड, जोड़ों में असंयुक्त (भिन्न जनकों की संतानें
अलग हैं क्योंकि जनक अलग थे; एक ही जनक की संतानें हटाए गए \omterm{prop:b1:reals:intervals}{अंतराल} से अलग हैं)।

\textbf{2.} प्रत्येक $C_n$ खंडों का परिमित सम्मिलन है, अतः \omterm{def:b1:topology:closed}{संवृत}; $C =
\bigcap C_n$ \omterm{def:b1:topology:closed}{संवृत} समुच्चयों का प्रतिच्छेदन है: अतः \omterm{def:b1:topology:closed}{संवृत}
(\cref{def:b1:topology:closed}); परिबद्ध ($\subseteq \intcc{0}{1}$): और
\cref{thm:b1:topology:compact} से संहत। अरिक्त: $0$ प्रत्येक $C_n$ के सबसे
बाएँ खंड में है। मान लीजिए $a$ $C_n$ के खंड $S$ का सिरा है। $m \leq n$ के
लिए $a \in C_n \subseteq C_m$। आगे के चरणों के लिए: मध्य-तिहाई हटाना कभी
किसी सिरे को नहीं हटाता, और $a$ फिर $S$ की दोनों संतानों में से एक का सिरा
है (वह संतान जो $a$ को छूती है); आगमन से सभी $m \geq n$ के लिए $a \in C_m$:
$a \in C$।

\textbf{3.} $C_n$ की कुल लंबाई: $2^n \cdot 3^{-n} = (\frac23)^n \to 0$।
$\varepsilon > 0$ दिया हो, तो ऐसा $n$ चुनिए कि $(\frac23)^n \leq
\varepsilon$: तब $C \subseteq C_n$, जो कुल लंबाई $\leq \varepsilon$ वाले
परिमित कितने खंडों का सम्मिलन है।

\textbf{4.} मान लीजिए $I \subseteq C$ दो भिन्न बिंदुओं वाला \omterm{prop:b1:reals:intervals}{अंतराल} है।
प्रत्येक $n$ के लिए: $I \subseteq C_n$, और $I$ उत्तल होने से $C_n$ के
\emph{एक ही} खंड के भीतर होना चाहिए (दो खंडों से मिलने पर $I$ को उनके बीच के
\omterm{prop:b1:reals:intervals}{अंतराल} का कोई बिंदु समेटना पड़ता, जो $C_n$ के बाहर है)। अतः $I$ की लंबाई सभी
$n$ के लिए $\leq 3^{-n}$ है: विरोधाभास। अतः $C$ के भीतर एकमात्र \omterm{prop:b1:reals:intervals}{अंतराल} रिक्त
या एकल हैं; विशेष रूप से कोई $\intoo{x-r}{x+r}$ $C$ के भीतर नहीं समाता:
$\mathring C = \emptyset$। और $C$ के \omterm{def:b1:topology:closed}{संवृत} होने से $\overline C = C$ का
\omterm{def:b1:topology:closure}{अंतःभाग} रिक्त है: $C$ कहीं \omterm{def:b1:topology:dense}{सघन} नहीं है।

\textbf{5.} $\varphi_0(x) = \frac x3$ और $\varphi_2(x) = \frac{2 + x}{3}$
लिखिए, जो $\intcc{0}{1}$ के $\intcc{0}{\frac13}$ तथा $\intcc{\frac23}{1}$ पर
वर्धमान ऐफ़ीन एकैकी आच्छादन हैं। दावा: $C_{n+1} = \varphi_0(C_n) \cup
\varphi_2(C_n)$। $n = 0$ के लिए यह प्रश्न 1 है। आगमन: वर्धमान ऐफ़ीन
\omterm{def:b1:logic:map}{प्रतिचित्रण} किसी खंड के मध्य-तिहाई को प्रतिबिंब खंड के मध्य-तिहाई पर भेजता
है, अतः मध्य-तिहाई हटाना $\varphi_0$ और $\varphi_2$ के साथ क्रम-विनिमेय है;
$C_{n+1} = \varphi_0(C_n) \cup \varphi_2(C_n)$ पर हटाने का पद लगाने से
$C_{n+2} = \varphi_0(C_{n+1}) \cup \varphi_2(C_{n+1})$ मिलता है। $n$ पर
प्रतिच्छेदन लेने पर: $x \leq \frac13$ के लिए सभी $n$ पर $x \in C \iff x \in
\varphi_0(C_n)$ $\iff 3x \in \bigcap C_n = C$; $\intcc{\frac23}{1}$ पर भी
वैसे ही; और $\intoo{\frac13}{\frac23}$ का कोई बिंदु $C_1$ में नहीं है। अतः
$C = \varphi_0(C) \cup \varphi_2(C)$, और वह भी असंयुक्त रूप से।

\textbf{6.} $n$ पर आगमन; स्थिति $n = 0$ कहती है कि प्रत्येक $x \in
\intcc{0}{1}$ का कोई कूट है, जो \cref{pb:b1:reals:1} है ($x < 1$ के लिए
प्रश्न 9; $1 = (0.\overline 2)_3$)। कोटि $n$ पर तुल्यता मान लीजिए। यदि $x
\in C_{n+1}$: तो प्रश्न 5 से $z \in C_n$ और $i \in \{0, 2\}$ के साथ $x =
\varphi_i(z)$; और यदि $(e_k)$ $z$ का ऐसा कूट है जिसके पहले $n$ अंक $1$-मुक्त
हैं, तो $(i, e_1, e_2, \dots)$ के आंशिक योग $\frac i3 + \frac13\sum_{k\leq
m} e_k 3^{-k} \to \varphi_i(z) = x$ हैं: अर्थात् $x$ का ऐसा कूट जिसके पहले
$n + 1$ अंक $1$-मुक्त हैं। विलोमतः, यदि $x$ का कोई कूट $(d_k)$ हो जिसमें
$d_1, \dots, d_{n+1} \in \{0, 2\}$: तो विस्थापित माला $(d_2, d_3, \dots)$ का
कोई मान $z \in \intcc{0}{1}$ है, उसके पहले $n$ अंक $1$-मुक्त हैं, और
आंशिक-योग की संगणना उलटी दिशा में पढ़ने पर $x = \varphi_{d_1}(z)$ देती है;
आगमन से $z \in C_n$, अतः प्रश्न 5 से $x \in C_{n+1}$। अंत में: पूर्णतः
$1$-मुक्त कूट $x$ को प्रत्येक $C_n$ में रख देता है, अतः $C$ में; विलोमतः यदि
$x \in C$, तो प्रत्येक $n$ के लिए $x$ के अधिक से अधिक दो कूटों
(\cref{pb:b1:reals:1}, प्रश्न 11) में से किसी एक के पहले $n$ अंक $1$-मुक्त
हैं; और कोई एक नियत कूट मनचाहे बड़े $n$ के लिए काम करना ही चाहिए (दो कूटों
के बीच कबूतरखाना), और जिस कूट के पहले $n$ अंक मनचाहे बड़े $n$ के लिए
$1$-मुक्त हों वह सीधे-सीधे $1$-मुक्त है।

\textbf{7.} $0 = (0.\overline 0)_3$, $1 = (0.\overline 2)_3$, $\frac13 =
(0.0\overline{2})_3$ ($(0.1)_3$ का अनुचित जुड़वाँ), $\frac23 =
(0.2\overline{0})_3$। गुणोत्तर योग:
\[
(0.\overline{02})_3 = \sum_{j\geq1} \frac{2}{9^{\,j}}
= \frac{2/9}{1 - 1/9} = \frac14 ,
\qquad
(0.\overline{20})_3 = \sum_{j\geq1} \frac{2}{3\cdot 9^{\,j-1}}
= \frac{2/3}{1 - 1/9} = \frac34 ,
\]
दोनों $1$-मुक्त: $\frac14, \frac34 \in C$। (अनंत योग आंशिक योगों के \omterm{def:b1:reals:bounds}{उच्चतम}
का संक्षेप हैं, जैसा \cref{pb:b1:reals:1} में है।) $C_n$ के खंडों के सिरों
का रूप $m/3^n$ है (आगमन: संतानों के सिरे जनक के सिरे हैं या उनमें से किसी से
$3^{-n-1}$ के गुणज भर भिन्न हैं)। यदि $\frac14 = \frac{m}{3^n}$, तो $3^n =
4m$, और $4 \nmid 3^n$: असंभव। अतः $\frac14 \in C$, और वह भी बिना कभी सिरा
हुए।

\textbf{8.} मान लीजिए $x$ के दो भिन्न $1$-मुक्त कूट होते। दो कूट होने का
अर्थ ही है (\cref{pb:b1:reals:1}, प्रश्न 11, आधार $3$) कि $x = m/3^N \in
\intoo{0}{1}$, और दोनों कूट ये हैं: समाप्त होने वाला कूट, जिसका अंतिम अशून्य
अंक $d_N \in \{1, 2\}$ है और उसके बाद $0$; और उसका जुड़वाँ, जिसमें स्थान $N$
पर $d_N - 1$ है और उसके बाद $2$। यदि $d_N = 1$, तो पहले में कोई $1$ है; और
यदि $d_N = 2$, तो जुड़वाँ में $d_N - 1 = 1$ है। दोनों में से किसी भी हाल में
उस जोड़े का अधिक से अधिक एक ही $1$-मुक्त है: विरोधाभास। अतः प्रत्येक $x \in
C$ का ठीक एक $1$-मुक्त कूट है (अस्तित्व प्रश्न 6 से), और भिन्न
$\{0,2\}$-मालाओं के मान भिन्न हैं। प्रत्येक $\{0,2\}$-माला का मान
$\intcc{0}{1}$ में है (आंशिक योग $\leq 1$) और उसके सभी उपसर्ग $1$-मुक्त हैं,
अतः मान प्रत्येक $C_n$ में, अर्थात् $C$ में है: मान-प्रतिचित्रण
$\{0,2\}$-मालाओं से $C$ पर एकैकी आच्छादन है।

\textbf{9.} मान लीजिए $(d^{(k)})$ $x_k$ का $1$-मुक्त कूट है और $e_k = 2 -
d^{(k)}_k \in \{0, 2\}$ रखिए: यह ऐसी $\{0,2\}$-माला है जिसका मान $y$ $C$ में
है और जिसका अद्वितीय $1$-मुक्त कूट $(e_k)$ है (प्रश्न 8)। प्रत्येक $k$ के
लिए $y$ और $x_k$ के कूट स्थान $k$ पर भिन्न हैं, अतः $y \neq x_k$: कोई भी
\omterm{def:b1:logic:map}{प्रतिचित्रण} $\N^* \to C$ \omterm{def:b1:logic:inj}{आच्छादक} नहीं है। लंबाई-शून्य वाला अगणनीय \omterm{def:b1:logic:sets}{समुच्चय}:
\omterm{def:b1:counting:card}{गणनांक} में विशालता और माप में लघुता --- एक साथ।

\textbf{10.} संहतता: अनंत प्रतिच्छेदन की संवृतता और परिबद्धता (प्रश्न 2) ---
यही एकमात्र गुणधर्म है जिसे कोई एक अंतर्विष्टि नहीं ढोती। लंबाई शून्य: $C
\subseteq C_n$, जिसकी कुल लंबाई $(\frac23)^n$ है (प्रश्न 3)। कहीं \omterm{def:b1:topology:dense}{सघन} न
होना: $C \subseteq C_n$ $C$ के भीतर के अंतरालों की लंबाई $\leq 3^{-n}$ होने
पर बाध्य कर देता है (प्रश्न 4)। अगणनीयता किसी अंतर्विष्टि पर सवार ही नहीं
है: उसे पूरी प्रतिच्छेदन-संरचना चाहिए, जो प्रश्न 8 के एकैकी आच्छादन में
कूटित है।

\textbf{11.} $d_n$ को $2 - d_n$ में पलटिए: नई माला अब भी $\{0,2\}$-माला है,
अतः उसका मान $x_n$ $C$ में है; और कोटि $n$ के आगे आंशिक योग ठीक
$2\cdot3^{-n}$ भर भिन्न हैं, अतः $\abs{x_n - x} = 2\cdot3^{-n}$। इस प्रकार
$x_n \neq x$ और $x_n \to x$: $C$ का प्रत्येक बिंदु $C$ के दूसरे बिंदुओं की
सीमा है --- अर्थात् $C$ \emph{पूर्ण} है, और उसका कोई वियुक्त बिंदु नहीं है।

\textbf{12.} प्रश्न 5 के आगमन से $C_n$ के खंड ठीक $\intcc{t}{t + 3^{-n}}$
हैं, जहाँ $t$ लंबाई-$n$ वाली $\{0,2\}$-मालाओं के मानों पर चलता है। कूट
$(d_k)$ वाला $x \in C$ दिया हो, तो कटाई $t_n$ (अंक $d_1 \dots d_n$, फिर $0$)
इसलिए एक बायाँ सिरा है, और $0 \leq x - t_n \leq 3^{-n}$: अतः सिरे $C$ में
\omterm{def:b1:topology:dense}{सघन} हैं। वे $\{m/3^n : m, n\}$ का उपसमुच्चय बनाते हैं, जो पूर्णांक युग्मों
से सूचीबद्ध है और इसलिए गणनीय है (जैसा \cref{exo:b1:topology:10} में $\Q$ के
लिए)। और चूँकि $C$ अगणनीय है (प्रश्न 9), $C$ के गणनीय कितने को छोड़कर सभी
बिंदु सिरे नहीं हैं --- $\frac14$ (प्रश्न 7) उसी हिमखंड की दिखाई देने वाली
नोक है।

\textbf{13.} ऐसा $n$ चुनिए कि $3^{-n} < y - x$। दोनों $x, y \in C_n$ हैं, और
वे एक ही खंड में नहीं हो सकते (लंबाई $3^{-n} < y - x$): उनके खंडों के बीच
हटाया गया \omterm{prop:b1:reals:intervals}{अंतराल} ऐसा $z$ देता है कि $x < z < y$ और $z \notin C_n \supseteq
C$। अतः $C$ के किन्हीं दो बिंदुओं को पूरक अलग कर देता है: $C$ के एकमात्र
उत्तल उपसमुच्चय एकल \omterm{def:b1:logic:sets}{समुच्चय} हैं --- अर्थात् $C$ पूर्णतः असंबद्ध है।

\textbf{14.} यदि $(d_k)$ $x$ का $1$-मुक्त कूट है, तो माला $(2 - d_k)$ फिर
$\{0,2\}$-माला है, जिसके आंशिक योग हैं
\[
\sum_{k=1}^{n} (2 - d_k)3^{-k} = (1 - 3^{-n}) - \sum_{k=1}^n
d_k 3^{-k} \longrightarrow 1 - x :
\]
अतः $1 - x \in C$। इस प्रकार $1 - C \subseteq C$, और इस \omterm{def:b1:logic:map}{प्रतिचित्रण} को दो
बार लगाने पर $1 - C = C$ मिलता है: अर्थात् \omterm{pb:b1:topology:1}{कैंटर समुच्चय} $\frac12$ के
सापेक्ष सममित है।

\textbf{15.} आंशिक योग $t_n \to x$ और $t'_n \to x'$ (वर्धमान अनुक्रम अपने
\omterm{def:b1:reals:bounds}{उच्चतम}, अर्थात् मान, की ओर अभिसरित होते हैं), अतः
\cref{thm:b1:seq:operations} से $t_n + t'_n \to x + x'$। कूट $(e_k)$ वाला $y
\in \intcc{0}{1}$ दिया हो, तो $e_k = 0, 1, 2$ के अनुसार $(a_k, b_k) = (0,0),
(0,2), (2,2)$ चुनिए: तब $a_k + b_k = 2e_k$, मालाएँ $(a_k)$, $(b_k)$
$\{0,2\}$-मालाएँ हैं जिनके मान $x, x' \in C$ हैं, और
\[
x + x' = \lim_n\,(t_n + t'_n) = \lim_n 2\sum_{k=1}^n e_k 3^{-k}
= 2y :
\]
अर्थात् प्रत्येक $y \in \intcc{0}{1}$ $C$ के दो बिंदुओं का मध्यबिंदु है।

\textbf{16.} प्रश्न 15 $\intcc{0}{2} = 2\,\intcc{0}{1} \subseteq C + C$ देता
है, और $C + C \subseteq \intcc{0}{1} + \intcc{0}{1} = \intcc{0}{2}$: अतः
समता। फिर $1 - C = C$ का प्रयोग करते हुए:
\[
C - C = C + (C - 1) = (C + C) - 1 = \intcc{-1}{1} .
\]
मूर्त उदाहरण: $1 = \frac14 + \frac34$, जो $C$ के दो गैर-सिरा सदस्यों का योग
है।

\textbf{17.} लंबाई यह नापती है कि \omterm{def:b1:logic:sets}{समुच्चय} स्वयं रेखा का कितना भाग घेरता है;
वह \emph{योगों} के \omterm{def:b1:logic:sets}{समुच्चय} के विषय में कुछ नहीं कहती, जो दो-प्राचल कुल $C
\times C$ का $(x, x') \mapsto x + x'$ के अंतर्गत प्रतिबिंब है --- दोनों
अंक-मालाएँ स्वतंत्र रूप से चुनी जाती हैं, और यही स्वतंत्रता $\intcc{0}{2}$
को भर देती है। कोई प्रमेय योग-समुच्चय की लंबाई को योज्यों की लंबाइयों से
परिबद्ध नहीं करती, और $C$ इस बात की उपपत्ति है कि कोई कर भी नहीं सकती।

\textbf{18.} \cref{pb:b1:reals:1} (प्रश्न 18) से $x$ परिमेय है तभी जब उसका
उचित प्रसार अंततः आवर्ती हो। $x \in C$ का $1$-मुक्त कूट या तो वही उचित
प्रसार है या किसी समाप्त होने वाले कूट का अनुचित जुड़वाँ; और समाप्त होने
वाली माला तथा उसका जुड़वाँ (अंततः अचर $2$) दोनों अंततः आवर्ती हैं, अतः
$1$-मुक्त कूट की आवर्तिता $x$ की परिमेयता के तुल्य है। आधार $3$ ($r_0 = 1$)
में $\frac1{13}$ का लंबा भाग: $3 = 13\cdot0 + 3$, $9 = 13\cdot0 + 9$, $27 =
13\cdot2 + 1$, और शेषफल $1$ पर लौट आता है: अंक $\overline{002}$, अतः
$\frac1{13} = (0.\overline{002})_3$, जो $1$-मुक्त और आवर्ती है: अर्थात् $C$
का एक परिमेय सदस्य। (जाँच: $\frac{2/27}{1 - 1/27} = \frac{2}{26} =
\frac1{13}$।)

\textbf{19.} जिस माला में त्रिभुजीय स्थानों $k = \frac{j(j+1)}{2}$ पर $d_k =
2$ हैं और अन्यत्र $0$, वह $\{0,2\}$-माला है, अतः उसका मान $x^*$ $C$ का सदस्य
है (प्रश्न 8)। उसमें अनंत कितने $2$ हैं जिनके क्रमागत \omterm{prop:b1:reals:intervals}{अंतराल} $j + 1 \to
\infty$ हैं, अतः वह अंततः आवर्ती नहीं है (आवर्त $T$ अंततः $\leq T$ अंतरालों
पर $2$ लाने को बाध्य करता: \cref{pb:b1:reals:1} (प्रश्न 20) का बढ़ते-अंतराल
वाला तर्क); और प्रश्न 18 से $x^* \notin \Q$। और प्रश्न 9 तथा $\Q$ की गणनीयता
से $C$ के गणनीय कितने को छोड़कर सभी सदस्य अपरिमेय हैं: $x^*$ नियम है, अपवाद
नहीं।

\textbf{20.} मान लीजिए $y \in \intcc{0}{1}$: उसका द्विआधारी कूट $(c_k)$ है,
जहाँ $c_k \in \{0, 1\}$ (\cref{pb:b1:reals:1}, प्रश्न 9, आधार $2$; $y = 1$
सर्व-$1$ वाली माला लेता है)। तब $(2c_k)$ $\{0,2\}$-माला है, उसका मान $x$ $C$
में है, और $h(x)$ $(c_k)$ का मान है, अर्थात् $y$: अतः $h$ $C$ को
$\intcc{0}{1}$ पर भेजता है। यदि $C$ $\N^*$ से किसी \omterm{def:b1:logic:map}{प्रतिचित्रण} का प्रतिबिंब
होता, तो $h$ के साथ संयोजन $\intco{0}{1}$ की पूरी सूची बना देता, जो
\cref{pb:b1:reals:1} (प्रश्न 22) की विकर्ण प्रमेय का विरोध करता: अतः $C$ फिर
से अगणनीय है। लंबाई-शून्य वाला \omterm{def:b1:logic:sets}{समुच्चय} पूरी लंबाई वाले खंड पर आच्छादित हो
रहा है।

\textbf{21.} प्रश्न 5 से $C$ $\varphi_0(C)$ और $\varphi_2(C)$ का असंयुक्त
सम्मिलन है, और हर एक मापित प्रति $\frac13 C$ का स्थानांतरण है। योगात्मकता,
मापन और स्थानांतरण-अपरिवर्त्यता देती हैं
\[
L(C) = L(\varphi_0(C)) + L(\varphi_2(C))
= \tfrac13 L(C) + \tfrac13 L(C) = \tfrac23\,L(C),
\]
अतः $\frac13 L(C) = 0$: $L(C) = 0$। केवल स्वसमरूपता ही $C$ को लंबाई-शून्य का
दंड सुना देती है --- प्रश्न 3 ने तो केवल दंड का पालन किया।

\textbf{22.} लंबाई $l_n$ का खंड $4^{-(n+1)}$ लंबाई का केंद्रीय \omterm{prop:b1:reals:intervals}{अंतराल} खोता
है, और $l_{n+1} = \frac{l_n - 4^{-(n+1)}}{2}$ लंबाई के दो खंड बच जाते हैं;
$l_0 = 1$ से आगमन $l_n = \frac{2^n + 1}{2\cdot4^n}$ की पुष्टि करता है:
वस्तुतः $\frac12\Bigl(\frac{2^n+1}{2\cdot4^n} - \frac{1}{4^{n+1}}\Bigr) =
\frac{2(2^n + 1) - 1}{2\cdot4^{n+1}} = \frac{2^{n+1} + 1}{2\cdot4^{n+1}}$,
और सदा $l_n > 0$: रचना कभी भूखी नहीं मरती। $K = \bigcap K_n$ \omterm{def:b1:topology:closed}{संवृत} और
परिबद्ध है, अतः संहत; और $K$ के भीतर का कोई \omterm{prop:b1:reals:intervals}{अंतराल} $K_n$ के एक ही खंड में
है, जिसकी लंबाई $l_n \to 0$ है: अर्थात् रिक्त \omterm{def:b1:topology:closure}{अंतःभाग}। हटाई गई लंबाई:
$\sum_{n\geq0} 2^n \cdot 4^{-(n+1)} =
\frac14\sum_{n\geq0}\bigl(\frac12\bigr)^n = \frac12$, और प्रत्येक $K_n$ की
कुल लंबाई $2^n l_n = \frac{2^n + 1}{2^{n+1}} > \frac12$ है। अब मान लीजिए
परिमित कितने \omterm{def:b1:topology:open}{विवृत} अंतरालों का सम्मिलन $U \supseteq K$ है।
\cref{ex:b1:topology:nested} के साथी \omterm{def:b1:logic:statement}{कथन} से किसी $n$ के लिए $U \supseteq
K_n$; और अंतरालों के परिमित सम्मिलनों पर लंबाई की योगात्मकता मान लेने पर
ढकने वाले अंतरालों की कुल लंबाई कम से कम $K_n$ जितनी है, जो $\frac12$ से
अधिक है। अतः $K$ कहीं \omterm{def:b1:topology:dense}{सघन} तो है, पर कोई सस्ता आवरण विद्यमान नहीं: सांस्थितिक
लघुता (कहीं \omterm{def:b1:topology:dense}{सघन} नहीं) और मापीय लघुता (लंबाई शून्य) सचमुच भिन्न धारणाएँ हैं,
और $K$ उन्हें अलग कर देता है।

\textbf{23.} प्राप्ति: मान लीजिए $d = d(x, F)$ और ऐसे $a_k \in F$ चुनिए कि
$\abs{x - a_k} \leq d + \frac1k$: $a_k$ परिबद्ध हैं, अतः
बोल्ज़ानो--वाइरश्ट्रास (\cref{thm:b1:seq:bw}) $a_{\varphi(k)} \to a$ निकाल
देता है, जहाँ $a \in F$ ($F$ \omterm{def:b1:topology:closed}{संवृत}, \cref{thm:b1:topology:seqclosed}) और
$\abs{x - a} = \lim \abs{x - a_{\varphi(k)}} = d$। अब अधिकतम: यदि $y \in C$,
तो $d(y, C) = 0$; अन्यथा $y$ किसी चरण $n \geq 1$ पर हटाए गए \omterm{prop:b1:reals:intervals}{अंतराल} में है,
जो लंबाई $3^{-n}$ का \omterm{def:b1:topology:open}{विवृत} \omterm{prop:b1:reals:intervals}{अंतराल} है और जिसके दोनों सिरे $C$ के सदस्य हैं
(प्रश्न 2), अतः $d(y, C) \leq \frac{3^{-n}}{2} \leq \frac16$, जहाँ समता के
लिए $n = 1$ और $y$ का उस \omterm{prop:b1:reals:intervals}{अंतराल} $\intoo{\frac13}{\frac23}$ के केंद्र पर होना
आवश्यक है, अर्थात् $y = \frac12$; और सचमुच $d\bigl(\frac12, C\bigr) =
\frac16$, क्योंकि $C \cap \intoo{\frac13}{\frac23} = \emptyset$ और $\frac13,
\frac23 \in C$। अतः $\max_{y\in\intcc{0}{1}} d(y, C) = \frac16$, जो ठीक
$\frac12$ पर प्राप्त होता है।

\textbf{24.} सिरे $C$ का गणनीय \omterm{def:b1:topology:dense}{सघन} उपसमुच्चय बनाते हैं (प्रश्न 12): उन्हें
एक ही अनुक्रम $(e_j)_{j\geq1}$ के रूप में सूचीबद्ध कीजिए, जो $C$ में एक
अनुक्रम है। उसकी सभी उपानुक्रमीय सीमाएँ $C$ में हैं ($C$ \omterm{def:b1:topology:closed}{संवृत})। विलोमतः, $x
\in C$ नियत कीजिए: प्रत्येक $n$ के लिए $x$ ($m \geq n$) को समेटने वाले $C_m$
के खंडों के सिरे $x$ के $3^{-m} \leq 3^{-n}$ के भीतर हैं, अतः अनंत कितने
भिन्न सिरे $x$ के $3^{-n}$ के भीतर हैं; ऐसे सूचकांक $j_1 < j_2 < \dots$
चुनिए कि $\abs{e_{j_n} - x} \leq 3^{-n}$: यह $x$ की ओर अभिसरित \omterm{def:b1:seq:subsequence}{उपानुक्रम} है।
अतः $(e_j)$ की उपानुक्रमीय सीमाओं का \omterm{def:b1:logic:sets}{समुच्चय} ठीक $C$ है --- अर्थात् एक ही
अनुक्रम अगणनीय कितने बिंदुओं पर गुच्छे बना रहा है, जो अभिसारी अनुक्रम का
विपरीत छोर है, क्योंकि उसका गुच्छा-समुच्चय एकल होता है।

\textbf{25.} (क) रचना ने ये खपाए: स्वेच्छ प्रतिच्छेदन के अंतर्गत \omterm{def:b1:topology:closed}{संवृत}
समुच्चयों का स्थायित्व ($C$ का \omterm{def:b1:topology:closed}{संवृत समुच्चय} के रूप में अस्तित्व), संहतता
प्रमेय \cref{thm:b1:topology:compact} (प्रश्न 2, 22, 23), और संवृतता तथा
अभिलग्नता के अनुक्रमीय अभिलक्षण (\cref{ex:b1:topology:nested} का
नेस्टेड-संहत तर्क और प्रश्न 23)। (ख) चारों युग्म: लंबाई शून्य फिर भी अगणनीय
(प्रश्न 3, 9); \omterm{def:b1:topology:closed}{संवृत} फिर भी रिक्त \omterm{def:b1:topology:closure}{अंतःभाग} वाला (प्रश्न 4); पूर्ण --- कोई
वियुक्त बिंदु नहीं --- फिर भी पूर्णतः असंबद्ध (प्रश्न 11, 13); नगण्य फिर भी
$C + C = \intcc{0}{2}$ वाला (प्रश्न 16)। (ग) प्रश्न 20 का आच्छादन $h$, संतत
और अह्रासमान बनाकर, संतत फलनों के सिद्धांत में शैतान की सीढ़ी बन जाता है; और
स्नातक वर्ष 3 के खंड के माप-सिद्धांत में $C$ वह मानक साक्षी है कि ``नगण्य''
का अर्थ ``गणनीय'' नहीं है, जबकि उसका मोटा चचेरा भाई (प्रश्न 22) ``कहीं \omterm{def:b1:topology:dense}{सघन}
नहीं'' को ``नगण्य'' से अलग कर देता है।
\end{solution}
